\section{Critérios operacionais e restrições de missão}

Para que o desenvolvimento do manipulador robótico seja estável, certos parâmetros foram definidos e serão apresentados a seguir.

\subsection{Conjunto de Juntas Admissíveis}
\label{sec:conjunto_juntas}

Conforme discutido na seção de Hipóteses de modelagem, as juntas do manipulador estão sujeitas a limites geométricos, cinemáticos e dinâmicos. Esses limites devem ser respeitados durante toda a operação de \emph{rendezvous and docking/berthing}.
Os limites permitem que os atuadores operem dentro de uma faixa segura estabelecida e que o manipulador não exceda as suas capacidades físicas.

Esses limites são formalizados pelos seguintes conjuntos admissíveis para posição articular \eqref{eq:conjunto_q_posicao}, velocidade \eqref{eq:conjunto_q_velocidade} e aceleração \eqref{eq:conjunto_q_aceleracao}:

\begin{equation}
\label{eq:conjunto_q_posicao}
\mathcal{Q}_{\text{pos}}^{adm} =
\left\{ q_j : q_{j,\min} \leq q_j \leq q_{j,\max} \right\}.
\end{equation}

\begin{equation}
\label{eq:conjunto_q_velocidade}
\mathcal{Q}_{\text{vel}}^{adm} =
\left\{ |\dot{q}_j| \leq \dot{q}_{j,\max} \right\}.
\end{equation}

\begin{equation}
\label{eq:conjunto_q_aceleracao}
\mathcal{Q}_{\text{acel}}^{adm} =
\left\{ |\ddot{q}_j| \leq \ddot{q}_{j,\max} \right\}.
\end{equation}

Na Tabela \ref{tab:limites_juntasPosVelocidade} são apresentados os limites\fn{Os limites da junta prismática, por ser linear, são dados em metro (m) e em metro/segundo (m/s)} de posição (em graus [°]) e velocidade (graus por segundo [°/s]) para cada uma das juntas. Os dados de aceleração máxima $\ddot{q}_{j,\max}$ serão apresentados posteriormente com base na análise dinâmica e no desempenho do controlador.

\begin{table}[ht]
\centering
\caption{Limites físicos e operacionais das juntas do manipulador robótico.}
\label{tab:limites_juntasPosVelocidade}
\begin{tabular}{clllll}
\toprule
\textbf{Junta} & \textbf{Tipo} & $q_{j,\min}$ & $q_{j,\max}$ & $\dot{q}_{j,\min}$ & $\dot{q}_{j,\max}$ \\
\midrule
1 & Revoluta & $-360^\circ$ & $360^\circ$ & $-90^\circ/\text{s}$ & $90^\circ/\text{s}$ \\
2 & Revoluta & $-90^\circ$ & $90^\circ$ & $-180^\circ/\text{s}$ & $180^\circ/\text{s}$ \\
3 & Revoluta & $-130^\circ$ & $130^\circ$ & $-180^\circ/\text{s}$ & $180^\circ/\text{s}$ \\
4 & Revoluta & $-150^\circ$ & $150^\circ$ & $-360^\circ/\text{s}$ & $360^\circ/\text{s}$ \\
5 & Prismática & $-0{,}15$ m & $0$ m & $-0{,}314$ m/s & $0{,}314$ m/s \\
\bottomrule
\end{tabular}
\end{table}

Os limites foram definidos de maneira incremental ao longo do desenvolvimento do manipulador robótico, de modo que possa refletir as capacidades dos atuadores, as restrições geométricas dos elos e as exigências da tarefa de \emph{berthing}.
Essa definição assegura a estabilidade do controle e previne a violação dos limites mecânicos durante a operação em microgravidade.




